\eabstract{
With quantitative trading increasingly dominating financial markets, trading strategies have formed a complex, interdependent ``strategy ecosystem.'' A core question of both theoretical and practical significance is whether strategy alpha (excess returns) systematically decays due to competition, imitation, and substitution among strategies. This study integrates three theoretical frameworks—Evolutionary Finance, the Adaptive Markets Hypothesis, and Agent-Based Computational Economics—to construct an Evolutionary Artificial Stock Market model (EVO-ASM) that systematically investigates alpha decay mechanisms in strategy ecosystems.

The model features 500 heterogeneous agents, a 6-dimensional signal space, Walrasian market clearing, and an evolution engine with elimination-replication-mutation operations. Through four progressive experiments—no-evolution baseline (E1), capital expansion (E2), replication mechanism (E3), and full evolution with replication and mutation (E4)—the causal mechanisms of alpha decay are disentangled, comprising several hundred independent simulation runs.

Experimental results demonstrate that: (1) alpha decay requires evolutionary mechanisms and does not occur systematically under static strategy distributions; (2) wealth concentration alone drives partial alpha decay, with each percentage point increase in crowding associated with a 0.3--0.8 basis point decline in alpha; (3) strategy replication accelerates diversity loss by a factor of 2--3; (4) mutation is critical for maintaining strategy ecosystem diversity, with moderate mutation rates establishing a dynamic equilibrium between competition and innovation; (5) the strategy ecosystem exhibits critical phase transitions, where the system may abruptly shift from a diverse to a homogeneous state—a finding with significant implications for financial stability monitoring. This study reconceptualizes alpha decay from an empirical residual to an emergent property of evolutionary dynamics, providing a new theoretical framework and analytical toolkit for understanding market efficiency dynamics in the quantitative era.
}{Evolutionary Finance; Adaptive Markets Hypothesis; Agent-Based Computational Economics; Artificial Stock Market; Alpha Decay; Strategy Ecosystem}
